\documentclass[aspectratio=1610]{beamer}
\usepackage[utf8]{inputenc}
\usepackage{booktabs}
\usepackage{tabulary}
\usepackage{amssymb}% http://ctan.org/pkg/amssymb
\usepackage{pifont}% http://ctan.org/pkg/pifont
\newcommand{\cmark}{\ding{51}}%
\newcommand{\xmark}{\ding{55}}%
\usetheme{simple}
\usecolortheme{whiteonblack}
\title{State of the Art in Linked Data Editors:\\Semantic Media Wiki and Alternatives}
%\subtitle{Semantic MediaWiki}
\author{Konrad Höffner}

\newcommand{\imageslide}[4][]
{
%\newgeometry{margin=0cm,top=1em}
\begin{frame}[plain]{~~~~#2}
\vspace{0.2em}
\centering\includegraphics[width=1.0\textwidth,height=0.95\textheight,keepaspectratio]{#3}
\\#1
\note{#4}
\end{frame}
%\restoregeometry
}

\begin{document}
\begin{frame}
\titlepage
\end{frame}

\begin{frame}{What is Linked (Open) Data?}
\centering
\includegraphics[height=0.5\textheight]{img/linked-data-mug.jpg}
\end{frame}

\imageslide{What is RDF}{img/spider.png}{}{}

\imageslide{Ontologies}{img/hitontology.pdf}{}{}

\imageslide{Linked Data Cloud}{img/lodcloud-excerpt-inverted.png}{}{}

\begin{frame}{Linked Data Lifecycle}
\centering
\includegraphics[height=0.9\textheight]{cycle.pdf}
\end{frame}

\begin{frame}{Storage---Virtuoso SPARQL Endpoint}
\centering
\includegraphics[width=\textwidth]{img/sparqlresult.png}
\end{frame}

\begin{frame}{Browse---LodView}
\centering
\includegraphics[width=\textwidth]{img/lodview.png}
\end{frame}

\begin{frame}{Explore---SNIK Graph}
\centering
\includegraphics[width=\textwidth]{img/snik-graph-full.png}
\end{frame}

\begin{frame}{Search---Facete 2}
\centering
\includegraphics[width=\textwidth]{img/facete2.png}
\end{frame}

\begin{frame}{Search---Facete 3}
\centering
\includegraphics[width=\textwidth]{img/facete3.png}
\end{frame}

\begin{frame}{Problem: How to add, delete and modify Linked Data?}
\centering
\includegraphics[height=0.9\textheight]{cycle-red.pdf}
\end{frame}

\begin{frame}[fragile]{Problem: How to add, delete and modify Linked Data?}
\begin{verbatim}
:HNA_Millenium rdf:type owl:NamedIndividual, :SoftwareProduct ;
 :offersFeature
   :allowing_to_reuse_earlier_anamnestic_data ,
   :automatic_creation_of_a_nursing_transfer_report ,
   :comprising_a_set_of_nursing_care_standards ,
   :nursing_anamnesis ,
   :nursing_care_planning_supported_by_pre-defined_nursing_care_standards ,
   :nursing_diagnosis_based_on_the_NANDA_taxonomy ,
   :reminding_and_documentation_of_nursing_tasks ,
   :scheduling ,
   :support_nursing_care_planning_for_certain_patient_cases ,
   :use_of_clinical_pathways ;
   rdfs:label "HNA Millenium"@en .
\end{verbatim}
\end{frame}



\begin{frame}{Native Tools}
\centering
\begin{tabulary}{\textwidth}{lLL}
\toprule
Tool		&Advantages\\					
\midrule
Text Editor	&quick, small diffs, version control\\
SPARUL		&many changes at once, immediately available\\
\bottomrule
\end{tabulary}
\end{frame}

\begin{frame}{Problem solved?}
\centering
\includegraphics[width=0.9\paperwidth]{img/texteditorsparul.png}
\end{frame}

\begin{frame}{No! Too technical.}
\centering
\includegraphics[width=0.9\paperwidth]{img/complicated.png}
\end{frame}

\iffalse
\begin{frame}{Problem: Authoring isn't easy}
\centering
\begin{tabulary}{\textwidth}{lLL}
\toprule
Tool		&Advantages						&Cons\\
\midrule
Text Editor	&speed, version control, 			&requires Turtle \& Git skills\\
SPARUL		&many changes at once 				&requires SPARUL syntax\\
Prot\'eg\`e	&\\
\bottomrule
\end{tabulary}
\end{frame}
\fi

\begin{frame}{What are we looking for?}
\centering
\textbf{Must-Have}
\begin{itemize}
\item graphical user interface
\item web application
\item add, delete and modify triples 
\item Virtuoso backend
\item intuitive 
\item fast
\end{itemize}
\textbf{Nice-To-Have}
\begin{itemize}
\item docker container 
\item forms and templates for instances 
\item version control
\item accounts
\item easy to setup and maintain
\item large community, positive outlook
\end{itemize}
\end{frame}

\begin{frame}{Relational Databases: PHPMyAdmin}
\centering
\includegraphics[width=0.9\paperwidth]{img/phpmyadmin.png}
\end{frame}

\begin{frame}{Relational Databases: Oracle Forms}
\centering
\includegraphics[width=0.7\paperwidth]{img/oracleforms.png}
\end{frame}

\imageslide{OntoWiki}{img/ontowiki.png}{}{}
\imageslide{OntoWiki Resource View}{img/ontowiki-ceo.png}{}{}
\imageslide{OntoWiki Editing}{img/ontowiki-ceo-edit.png}{}{}

\begin{frame}{OntoWiki Conclusions}
\centering
\begin{itemize}
\item user management
\item history
\item docker container
\item multiple knowledge bases
\item well suited for editing instances 
\item not maintained anymore
\item slow on large amounts (millions) of triples
\end{itemize}
\end{frame}

\imageslide{Prot\'eg\'e, WebProt\'eg\'e}{img/protege.png}{}{}

\begin{frame}{Prot\'eg\'e, WebProt\'eg\'e Conclusions}
\centering
\includegraphics[width=0.4\paperwidth]{img/protege.png}
\begin{itemize}
\item reasoning
\item active community
\item cannot use Virtuoso backend
\item useful for ontology creation, not instances
\end{itemize}
\end{frame}

\begin{frame}{Semantic MediaWiki}
\centering
\includegraphics[width=0.6\paperwidth]{img/smw-tokyo.png}
\begin{enumerate}
\item what is the core idea?
\item can we use it to edit RDF?
\item is it user friendly? does it have enough functionality?
\item can it be integrated into our service infrastructure?
\item what is the environment, community and outlook?
\end{enumerate}
\end{frame}

\begin{frame}{Semantic MediaWiki---Core Idea}
\begin{itemize}
\item extension to MediaWiki (Wikipedia)
\item "store and query data within the wiki's page" 
\item hybrid between document and data
\item great fit for documents with data, but does it make sense for just data?
\end{itemize}
\end{frame}

\begin{frame}{Semantic MediaWiki---Edit RDF}
\begin{itemize}
\item typed links between documents
\item can be translated to RDF
\item PageForms extension to generate web forms with templates for a category $\rightarrow$ user generation of instances
\end{itemize}
\end{frame}

\begin{frame}[fragile]{Semantic MediaWiki---Usability without Forms}
\begin{verbatim}
{{#ask:
[[Kategorie:Stadt]]
[[Land::Schweiz]]
|?Name
|?Einwohner
}}
\end{verbatim}
\begin{itemize}
\item you need to check out for yourself whether it is usable for you
\item users need to learn Semantic MediaWiki syntax
\item we have a test account at an existing installation thanks to BITPlan
\item ask Konrad for credentials
\item be mindful, it is not our own instance
\end{itemize}
\end{frame}

\imageslide{Semantic MediaWiki---Create Form}{img/smw-create-form.png}{}{}
\imageslide{Semantic MediaWiki---Create Class}{img/smw-create-class.png}{}{}
%\imageslide{Semantic MediaWiki---}{img/smw-create-product-error.png}{}{}

\begin{frame}{Semantic MediaWiki---Problems}
\centering
\includegraphics[width=0.7\paperwidth]{img/smw-create-product-error.png}
\begin{itemize}
\item forms do not work with links to external knowledge bases
\item forms are static, need to be manually synchronized to  ontology changes
\item but all tested containers were broken
\item there is no out-of-the box option to use Virtuso as a backend
\end{itemize}
\end{frame}

\begin{frame}{Semantic MediaWiki---Community and Outlook}
\begin{itemize}
\item thanks to Wolfgang Fahl of BITPlan for an introduction
\item 200-300 MediaWiki developers, focus on MediaWiki, rest of the world has to follow them
\item Semantic MediaWiki lacks unified leadership, many different fractions
\item usual data store are SQL databases and to a lesser extend wikiBase
\item very few users of SPARQL endpoints, no public example 
\end{itemize}
\end{frame}

\begin{frame}{Semantic MediaWiki---Summary}
\centering
\begin{tabulary}{\textwidth}{lL}
\toprule
Property		&Assessment\\
\midrule
Edit RDF		&\cmark\\
Usability		&your assessment\\
Integration		&\xmark\\
Outlook			&unclear\\
\midrule
Fit for our purposes	&not quite\\
\bottomrule
\end{tabulary}
\end{frame}

\imageslide{InstanceGenerator---HITO Prototype}{img/instancegenerator.png}{}{}

\begin{frame}{Conclusions and Recommendation}
\begin{itemize}
\item there is no perfect solution, lack of tool projects
\item SemanticMediaWiki not a good fit for HITO in my opinion
\item keep using the OntoWiki for SNIK
\item put some time into fixing the OntoWiki connection to the newest Virtuoso version
\item discuss OntoWiki for HITO
\item use Prot\'eg\`e for ontologies like the HITO ontology
\item don't dismiss the manual approach! Git + text editor + SPARUL
\item use the instancegenerator for appropriate tasks
\end{itemize}
\end{frame}

\end{document}

